\subsection{Explicit global growth from certified heights}
\label{sec:fd-global-growth-numeric}

The arithmetic certificate also gives explicit leading constants
in the global detection laws. We first control the multiplicity at an
individual ordinate by an elementary positive logarithmic derivative.

\begin{lemma}[An explicit multiplicity bound]
\label{lem:fd-multiplicity-log}
Every nontrivial zero at positive ordinate \(\gamma\ge10\) has
multiplicity
\[
 m_\gamma<\frac34\log\gamma.
\]
\end{lemma}
\begin{proof}
The paired Hadamard product and the zero-location theorem give
\[
 \frac{2m_\gamma}{3}
 \le\Re\frac{\xi'}{\xi}(2+i\gamma).
\]
Indeed, the zero at \(1/2+i\gamma\) contributes exactly
\(m_\gamma/(3/2)\), and every other real part in the paired
logarithmic-derivative sum is positive.

For real \(y\), differentiating Euler's product for the gamma
function gives
\[
 \Re\psi(1+iy)=-\gamma_E+
       \sum_{n\ge1}\frac{y^2}{n(n^2+y^2)}.
\]
The summand, as a function of positive real \(n\), decreases.
Bounding its sum by its first term plus its integral from one to
infinity yields
\[
 \Re\psi(1+iy)
 \le-\gamma_E+\frac{y^2}{1+y^2}
                +\frac12\log(1+y^2)
 <1-\gamma_E+\frac12\log(1+y^2).
\]
The absolutely convergent Euler series gives
\[
 A_2=-\frac{\zeta'}{\zeta}(2)=\sum_{n\ge2}\frac{\Lambda(n)}{n^2}
 <\frac35.
\]
Here this last rational bound has a short finite certificate: summing
the \(193\) prime-power terms through \(1000\) and using
\(\Lambda(n)\le\log n\) gives
\[
 A_2\le\sum_{n\le1000}\frac{\Lambda(n)}{n^2}
             +\frac{\log1000+1}{1000}
 <0.57687<\frac35.
\]
The same scalar certificate verifies
\(\gamma_E\ge\sum_{j=1}^{6}j^{-1}-\log7>1/2\) and
\(\log(2\pi)>9/5\).

The completion formula, with \(y=\gamma/2\), now implies
\begin{align*}
 \Re\frac{\xi'}{\xi}(2+i\gamma)
 &<\frac3{\gamma^2}+\frac35-\frac12\log\pi
     +\frac12-\frac12\gamma_E
     +\frac14\log(1+\gamma^2/4)\\
 &=\frac12\log\gamma+\frac3{\gamma^2}+\frac{11}{10}
     -\frac12\{\gamma_E+\log(2\pi)\}
     +\frac14\log(1+4/\gamma^2)\\
 &<\frac12\log\gamma+\frac4{\gamma^2}-\frac1{20}
 <\frac12\log\gamma\qquad(\gamma\ge10).
\end{align*}
In the first line, the rational completion terms are bounded by
\(2/\gamma^2+1/\gamma^2\), and the Euler series has real part at
most \(A_2\). The final two inequalities use
\(\log(1+x)<x\) and \(4/\gamma^2\le1/25\).
Combining the first and last bounds proves the lemma.
\end{proof}

Write the extended global limits as
\[
 L_T=\lim_{a\downarrow0}\sqrt a\,m_*(-a)
     =\lim_{a\downarrow0}\sqrt a\,R_*(-a),\qquad
 L_J=\lim_{a\downarrow0}a\,j_*(-a).
\]

\begin{corollary}[Certified global constants and their comparison]
\label{cor:fd-global-growth-numeric}
For the actual zeta heat family,
\[
 L_T>1.65\cdot10^{12},\qquad L_J>4.5\cdot10^{24}.
\]
The arithmetic trace certificate through $97.632$ by itself gives
$L_T>216.88$ and $L_J>4766.003712$. The larger bounds use the
Platt--Trudgian simplicity certificate through $3\cdot10^{12}$
\cite[abstract, p.~792; preprint Section~2, pp.~2--3]{platt-trudgian2021}.
Moreover,
\[
 \frac{4\pi}{3}
     \sqrt{\frac{L_J}{\log\sqrt{2L_J}}}
 \le L_T\le
 \frac{\sqrt3\,\pi}{2}
     \sqrt{L_J\log\sqrt{2L_J}}.
\]
The two functions of \(L_J\) are interpreted as \(+\infty\) at
\(+\infty\).
\end{corollary}
\begin{proof}
The arithmetic certificate proves that the first multiple ordinate
\(\gamma_*\) exceeds \(97.632\), including the value infinity.
The exact limits in Theorem~\ref{thm:fd-limits}, with \(x=2\gamma\),
are
\[
 L_T=\min_{m_\gamma\ge2}\frac{\gamma\mathfrak r_{m_\gamma}}{\sqrt2},
 \qquad L_J=\frac{\gamma_*^2}{2}.
\]
Thus \(L_J>(97.632)^2/2=4766.003712\).

For multiplicities two and three, the exact radii
\(\mathfrak r_2=\pi\), \(\mathfrak r_3=2\pi/\sqrt3\) give a
Toeplitz contribution strictly greater than
\[
 \frac{97.632\,\pi}{\sqrt2}>216.88.
\]
For an integer multiplicity \(m\ge4\), the preceding lemma gives
\(\gamma>e^{4m/3}\). Proposition~\ref{prop:orbital-unimodal-bound}
therefore gives
\[
 \frac{\gamma\mathfrak r_m}{\sqrt2}
 >\pi\sqrt{\frac23}\,e^{4m/3}\frac{\sqrt{m+1}}m.
\]
The logarithmic derivative of the last function of real \(m\) is
\(4/3+1/(2(m+1))-1/m>0\) for \(m\ge4\). Its value at four is
\[
 \frac\pi4\sqrt{\frac{10}3}\,e^{16/3}>297.
\]
This covers every possible multiple zero and proves the stated
lower bound for \(L_T\).

For the comparison, the same orbital estimate implies
\(\mathfrak r_m>2\pi/\sqrt{3m}\). At every multiple ordinate,
Lemma~\ref{lem:fd-multiplicity-log} consequently gives
\[
 \frac{\gamma\mathfrak r_m}{\sqrt2}
 >\frac{2\sqrt2\,\pi}{3}
             \frac{\gamma}{\sqrt{\log\gamma}}.
\]
The function \(\gamma/\sqrt{\log\gamma}\) increases for
\(\gamma\ge10\), so the minimum is bounded below by its value
at \(\gamma_*\). For the upper bound, evaluate the same minimum
at the first multiple ordinate and use the previously proved
\(\mathfrak r_m\le\pi\sqrt{m-1}\). This gives
\[
 L_T\le\pi\sqrt{\frac38}\,
          \gamma_*\sqrt{\log\gamma_*}.
\]
Substituting \(\gamma_*=\sqrt{2L_J}\) proves both displayed
comparisons for finite \(\gamma_*\). The infinite case follows
from the exact simplicity equivalences in
Theorem~\ref{thm:fd-limits} and the stated extended interpretation.

Finally, the Platt--Trudgian certificate gives
\(\gamma_*>H:=3\cdot10^{12}\). Substitution into the exact
Jensen limit and the increasing lower comparison just proved gives
\[
 L_J>H^2/2=4.5\cdot10^{24},\qquad
 L_T>\frac{2\sqrt2\pi}{3}\frac{H}{\sqrt{\log H}}
                         >1.65\cdot10^{12}.
\]
For the last strict comparison it suffices to use
\(\sqrt2>1.414\), \(\pi>3.14\), \(\log H<28.8\), and
\(\sqrt{28.8}<5.367\); their substitution gives
\(2(1.414)(3.14)/5.367>1.65\). These inequalities are also
certified by exact rational and outward-ball arithmetic.
\end{proof}

All decimal comparisons in this subsection are strict outward-ball
inequalities in the accompanying scalar certificate. The smaller input height
is the arithmetic trace certificate, the larger is the cited
Platt--Trudgian packet, and the Euler sum
uses prime powers only through \(1000\). These are fixed lower bounds
and a quantitative comparison of the global constants. The question
whether the constants are infinite remains the all-height simplicity
question.
